\section{Erd\H{o}s Problem 178: one colouring for countably many ordered sets}

\subsection{1) Statement and scope}
Let $A_i=\{a_{i1}<a_{i2}<\cdots\}\subseteq\mathbb N$ be infinite,
for $i\geq1$. Is there a single $f:\mathbb N\to\{-1,1\}$ and constants
$C(d)$ depending only on $d$ such that
\[
 \left|\sum_{j=1}^m f(a_{ij})\right|\leq C(d)
 \quad(m\geq1,\ 1\leq i\leq d)
\]
for every $d$? Beck proved the affirmative answer. The elementary
constructions below solve each fixed finite-family problem and give
the sharp bound one for two sets. They do not produce the required
single colouring for the entire countable family.

\subsection{2) A finite family: partition by membership pattern}
Fix $d$. For $S\subseteq[d]$, put
\[
 B_S=\{n\in\mathbb N:\{i\leq d:n\in A_i\}=S\}.
\]
Within each nonempty $B_S$, list its elements in increasing order and
give them alternating signs, starting with $+1$. This works for both
finite and infinite classes. Every initial segment of a class has sum
zero or one.

The first $m$ elements of $A_i$ are precisely its elements at most
$x=a_{im}$. For each $S$ containing $i$, their intersection with $B_S$
is an initial segment of that class. Since there are $2^{d-1}$ such
patterns, adding their sums proves
\[
 \left|\sum_{j=1}^m f(a_{ij})\right|\leq2^{d-1}
 \quad(i\leq d,\ m\geq1).                              \tag{1}
\]
This is a uniform bound over all infinite sets in any fixed finite family.

\subsection{3) A sharper construction for two sets}
For each of $A_1,A_2$, pair its first and second elements, third and
fourth elements, and so on. Regard these pairs as edges of a graph on
$\mathbb N$. Each of the two edge collections is a matching. Their
union has maximum degree two, and every cycle has even length because
successive edges alternate between the two matchings. An edge present
in both matchings cannot belong to a longer cycle.

Consequently every component is bipartite. To specify a colouring
without a choice ambiguity, give the least vertex in a component sign
$+1$, and give any other vertex the sign determined by the parity of
a path from that least vertex. Even cycles make this independent of
the path. Every designated pair now has sum zero. A prefix of either
$A_i$ consists of complete pairs and at most one additional element,
so its sum has absolute value at most one. The first element of any
set already has absolute sum one, so this bound is best possible.

\subsection{4) Why this does not finish the countable problem}
The colouring in (1) depends on the chosen number $d$. If we apply it
to the first $D$ sets and let $D$ increase, the available bound for
$A_1$ is $2^{D-1}$, which also increases. A subsequential limit of
these colourings therefore need not have any finite discrepancy bound
even on $A_1$. Compactness would apply if bounds for the first $d$ sets
were uniform in all later $D$, but that is exactly the missing control.
The pairing construction also does not extend directly: three matchings
can contain an odd cycle.

\textbf{Outcome of this attempt: UNRESOLVED for the countable family.}
The finite-family exponential bound and the two-set bound are fully
proved. The established theorem of Beck has a stronger simultaneous
conclusion; its proof is not reconstructed here.
\begin{itemize}
\item J. Beck, the 1981 resolution cited as [Be81] on the problem page,
and the later 2017 quantitative improvement cited there as [Be17].
\item T. F. Bloom, Problem 178,
\texttt{https://www.erdosproblems.com/178}, checked 23 September 2026.
\end{itemize}
The percentage records partial proof coverage, not a claim that the
known problem itself lacks a solution.
\subsection{5) COMPLETION ESTIMATE}
\textbf{COMPLETION: 20\%}
